\chapter{Описание проекта}
\label{ch:project}

\section{Технологический стек}

Все четыре урока решаются исключительно средствами браузера и
\texttt{Node.js} --- без сборщиков, транспиляторов и менеджеров
зависимостей. Это сознательное решение: цель курса --- показать,
насколько мало <<ритуала>> требуется современному фронтенду.

\begin{itemize}
  \item \textbf{HTML5} --- разметка единственной страницы каталога;
  \item \textbf{CSS3 + Bootstrap~5.0.1} --- готовая 12-колоночная
        сетка, утилитарные классы, компоненты
        \texttt{card}, \texttt{modal}, \texttt{navbar},
        \texttt{spinner};
  \item \textbf{Vanilla JavaScript} (ES2022) --- никаких фреймворков
        и сборки; синтаксис \texttt{async/await},
        \texttt{template literals}, \texttt{Array.map}, опциональная
        цепочка \texttt{?.}; код выполняется браузером как есть;
  \item \textbf{Node.js~18+} --- запуск консольных задач уроков~1
        и~2 без браузера.
\end{itemize}

Bootstrap подключён одной строкой в \texttt{<head>}, JS-бандл
Bootstrap (для модальных окон) --- в конце \texttt{<body>}; никаких
\texttt{npm install} и \texttt{node\_modules}.

\section{Учебный API персонажей Marvel}

Источник данных --- учебный REST-API курса:
\url{https://jsfree-les-3-api.onrender.com/characters}. Без
авторизации, HTTPS, CORS открыт. По GET-запросу возвращает
JSON-массив из девятнадцати героев. Структура одного объекта:

\begin{lstlisting}[caption={Фрагмент ответа API (один персонаж)},captionpos=b]
{
    "id": 1009610,
    "name": "Человек-паук",
    "nameor": "Spider-Man",
    "description": "Укушенный радиоактивным пауком...",
    "modified": "2020-07-21T10:30:10-0400",
    "thumbnail": "http://i.annihil.us/u/prod/marvel/i/mg/3/50/526548a343e4b.jpg",
    "comics": [
        { "resourceURI": "...", "name": "A+X (2012) #4" },
        ...
    ],
    "urls": [ ... ]
}
\end{lstlisting}

Полезные поля: \texttt{id} (уникальный идентификатор --- используется
для связи карточки с модальным окном), \texttt{name} (русское имя),
\texttt{nameor} (оригинальное имя), \texttt{description} (биография),
\texttt{thumbnail} (URL обложки), \texttt{comics[]} (список комиксов).
Поле \texttt{modified} и \texttt{urls} в проекте не используются.

\textbf{Особенность.} Изображения раздаются по протоколу \texttt{http},
а сайт публикуется на \texttt{https}. Браузер блокирует такой
\textit{mixed content}, поэтому в \texttt{index.js} URL картинки
принудительно переписывается на \texttt{https} --- домен
\texttt{i.annihil.us} поддерживает оба протокола.

API хостится на бесплатном тарифе платформы Render. При отсутствии
запросов более 15 минут контейнер усыпляется и первый запрос может
обрабатываться до \(30\) секунд (\textit{cold start}). Это учтено в
интерфейсе: на время загрузки в \texttt{\#character-card-box} отображается
спиннер Bootstrap, ошибка сети выводится понятным сообщением в
\texttt{alert-danger}.

\section{Структура каталога \texttt{lab13/}}

Папка лабораторной работы разделена на четыре подкаталога по числу
уроков:

\begin{lstlisting}[caption={Файловая структура \texttt{lab13/}},captionpos=b]
lab13/
├── README.md
├── Методичка (lab13).pdf
├── les-1/
│   ├── task1.js          # пользователь + бонусный баланс
│   └── task2.js          # баланс через 7 дней
├── les-2/
│   ├── task1.js          # вывод переписки
│   └── task2.js          # поиск по includes
├── les-3/                # сайт «Персонажи Marvel»
│   ├── index.html
│   ├── index.js
│   └── start.js
└── latex-report/         # этот отчёт
\end{lstlisting}

\section{Архитектура сайта}

Сайт \texttt{les-3/} --- это \emph{single-page application} (SPA),
собирающий разметку прямо в браузере. Точка входа ---
\texttt{index.html}: статический каркас на Bootstrap-компонентах,
содержит навбар, заголовок, два пустых контейнера
(\texttt{\#character-card-box} и \texttt{\#character-modal-box}) и
спиннер-плейсхолдер. По событию \texttt{DOMContentLoaded}
вызывается функция \texttt{start()} из \texttt{start.js}.

\begin{lstlisting}[caption={Граф зависимостей файлов},captionpos=b]
index.html
    │
    ├── Bootstrap CSS (CDN, <head>)
    ├── Bootstrap JS bundle (CDN, конец <body>)
    │
    ├── index.js  ← fetchCharacters, getCharacterCard[s],
    │                getCharacterModal[s], escapeHtml, toHttps
    │
    └── start.js  ← start(): fetchCharacters → render
\end{lstlisting}

Разделение на \texttt{index.js} (чистые функции преобразования данных
в HTML-строки) и \texttt{start.js} (оркестрация: запрос, рендер,
обработка ошибок) повторяет учебный шаблон преподавателя и облегчает
тестирование отдельных функций.

\section{Способ публикации}

Публикация сайта (урок~4 методички) выполнена не на коммерческом
хостинге \textit{Timeweb}, а на собственном домашнем сервере под
управлением \texttt{Debian~12}. Сервер хостит несколько Docker-сервисов
(Crafty, FileBrowser, FindMyDevice) под одним обратным прокси
\textbf{Caddy}. Marvel-сайт добавляется в эту инфраструктуру как ещё
один Docker-контейнер на образе \texttt{nginx:alpine}, раздающий
статические файлы \texttt{index.html}, \texttt{index.js},
\texttt{start.js}.

Логическая схема:

\begin{lstlisting}[caption={Маршрут запроса до сайта},captionpos=b]
браузер  →  443/tcp  →  Caddy (Let's Encrypt TLS)
                              │
                              └─→ marvel:80 (nginx) → /usr/share/nginx/html
\end{lstlisting}

Адрес сайта --- \url{https://marvel.server34.netcraze.club}. TLS-сертификат
выпускается автоматически Let's~Encrypt при первом запросе по новому
поддомену. DNS-запись \texttt{marvel} --- CNAME на основной домен
\texttt{server34.netcraze.club}, который указывает на внешний IP
домашнего роутера; роутер пробрасывает порты \(80\) и \(443\) на сервер.

Автоматизация деплоя --- идемпотентный bash-скрипт
\texttt{home-server/scripts/51-deploy-marvel-site.sh}: синхронизирует
файлы сайта в \texttt{/opt/stack/marvel/site/}, добавляет (если ещё нет)
сервис \texttt{marvel} в общий \texttt{docker-compose.yml}, дописывает
блок поддомена в \texttt{Caddyfile} и перезагружает Caddy. Подробное
описание шагов скрипта --- в~\hyperref[ch:practice]{главе~3}.

\endinput
